\scanpage{45}
\margin{44}
\mspage{-3.6-}

for every $j\in\{0,1,\ldots,n\}$.

Note that $\lambda_j$ is a section of $L_j\otimes L_{\mu(1)}^{-1}$ over the open set $U_R$. We now define
\[
\phi:U_R\longrightarrow\mathrm{PGL}(n)\times\mathbb{A}_R
\]
by its two components $p_1\circ\phi$ and $p_2\circ\phi$. The projective group $\mathrm{PGL}(n)$, is, by definition, the affine open subset of
\[
\mathbb{P}_{n^2+2n}=\operatorname{Proj}\mathbb{Z}[a_{00},\ldots,a_{0n};\ a_{10},\ldots,a_{1n};\ \ldots;\ a_{n0},\ldots,a_{nn}]
\]
where $\det(a_{ij})\ne0$. Therefore $p_1\circ\phi$ is determined by
\[
\begin{aligned}
(p_1\circ\phi)^*(\underline{o}_{\mathbb{P}}(1))&=L_{\mu(1)}\\
(p_1\circ\phi)^*(a_{ij})&=x_i^{(j)}\otimes(-1)^j\lambda_j^{-1}.
\end{aligned}
\]

Finally the morphism $p_2\circ\phi$ into affine space is determined by:
\[
(p_2\circ\unsure{\psi|\phi})^*(x_i^{(n+k)})
=\frac{D_{0,1,\ldots,\widehat{i},\ldots,n,n+k}}{D_{0,1,\ldots,\widehat{\mu(k)},\ldots,n,n+k}}
\cdot\frac{\lambda_i}{\lambda_{\mu(k)}}
\]
(It should be noted that the expression on the right is actually a section of $\underline{o}_{\mathbb{A}_R}$). Modulo the mechanical verifications that remain, this proves the Proposition. \uline{QED}

Of the morphisms constructed in the course of the above proof, it is useful to \deleted{\unsure{let}} fix the notation:
\[
\phi_R=p_2\circ\phi:U_R\longrightarrow\mathbb{A}^{n\nu-n}
\]
